% 版式母版：A4 双面教材，正文和栏目各有明确的页面职责。
% 本文件是版式的唯一来源；章节文件只调用这里定义的接口。

\usepackage{geometry}
\geometry{
  a4paper,
  inner=22mm,
  outer=27mm,
  top=21mm,
  bottom=25mm,
  headheight=15pt,
  headsep=7mm,
  footskip=11mm
}

\usepackage{amsmath,mathtools}
% 数学字形也随稿固定；STIX Two Math 覆盖常用符号、字母和可伸缩定界符。
% 这样公式不再随发行版的 Computer Modern 版本变化。
\usepackage{unicode-math}
\setmathfont[
  Path=fonts/,
  Extension=.otf
]{STIXTwoMath-Regular}
\usepackage{booktabs,array,tabularx,longtable,makecell}
\usepackage{graphicx}
\usepackage{wrapfig}
\usepackage{enumitem}
\usepackage{xcolor}
\usepackage{tikz}
\usetikzlibrary{calc,positioning,arrows.meta,fit}
\usepackage[most]{tcolorbox}
\usepackage{fancyhdr}
\usepackage{titlesec}
\usepackage{titletoc}
\usepackage{setspace}
\usepackage{needspace}
\usepackage{etoolbox}
\usepackage{xparse}
\usepackage{hyperref}
\usepackage{bookmark}

% 正文字体固定为仓库内的 Noto CJK 文件；不读取操作系统字体，也不依赖
% ctex 在 Windows、Linux 或 macOS 上选择不同的 fontset。字体文件均保持
% 上游文件名和字形内容，版本与授权记录在 fonts/README.md 和 licenses/。
\defaultfontfeatures{Ligatures=TeX}
\setmainfont[
  Path=fonts/,
  Extension=.otf,
  UprightFont=NotoSerifSC-Regular,
  BoldFont=NotoSerifSC-Bold,
  ItalicFont=NotoSerifSC-Regular,
  BoldItalicFont=NotoSerifSC-Bold
]{NotoSerifSC-Regular}
\setsansfont[
  Path=fonts/,
  Extension=.otf,
  UprightFont=NotoSansCJKsc-Regular,
  BoldFont=NotoSansCJKsc-Bold,
  ItalicFont=NotoSansCJKsc-Regular,
  BoldItalicFont=NotoSansCJKsc-Bold
]{NotoSansCJKsc-Regular}
\setmonofont[
  Path=fonts/,
  Extension=.otf,
  UprightFont=NotoSansMonoCJKsc-Regular,
  BoldFont=NotoSansMonoCJKsc-Bold,
  ItalicFont=NotoSansMonoCJKsc-Regular,
  BoldItalicFont=NotoSansMonoCJKsc-Bold
]{NotoSansMonoCJKsc-Regular}
\setCJKmainfont[
  Path=fonts/,
  Extension=.otf,
  UprightFont=NotoSerifSC-Regular,
  BoldFont=NotoSerifSC-Bold,
  ItalicFont=NotoSerifSC-Regular,
  BoldItalicFont=NotoSerifSC-Bold
]{NotoSerifSC-Regular}
\setCJKsansfont[
  Path=fonts/,
  Extension=.otf,
  UprightFont=NotoSansCJKsc-Regular,
  BoldFont=NotoSansCJKsc-Bold,
  ItalicFont=NotoSansCJKsc-Regular,
  BoldItalicFont=NotoSansCJKsc-Bold
]{NotoSansCJKsc-Regular}
\setCJKmonofont[
  Path=fonts/,
  Extension=.otf,
  UprightFont=NotoSansMonoCJKsc-Regular,
  BoldFont=NotoSansMonoCJKsc-Bold,
  ItalicFont=NotoSansMonoCJKsc-Regular,
  BoldItalicFont=NotoSansMonoCJKsc-Bold
]{NotoSansMonoCJKsc-Regular}
% ctex 的宋、黑、仿宋、楷书接口都指向同一组本地字形，避免样式命令
% 在某个平台重新触发 SimSun、KaiTi 或 Fandol 的系统回退。
\setCJKfamilyfont{zhsong}[
  Path=fonts/,
  Extension=.otf,
  UprightFont=NotoSerifSC-Regular,
  BoldFont=NotoSerifSC-Bold,
  ItalicFont=NotoSerifSC-Regular,
  BoldItalicFont=NotoSerifSC-Bold
]{NotoSerifSC-Regular}
\setCJKfamilyfont{zhhei}[
  Path=fonts/,
  Extension=.otf,
  UprightFont=NotoSansCJKsc-Regular,
  BoldFont=NotoSansCJKsc-Bold,
  ItalicFont=NotoSansCJKsc-Regular,
  BoldItalicFont=NotoSansCJKsc-Bold
]{NotoSansCJKsc-Regular}
\setCJKfamilyfont{zhyahei}[
  Path=fonts/,
  Extension=.otf,
  UprightFont=NotoSansCJKsc-Regular,
  BoldFont=NotoSansCJKsc-Bold,
  ItalicFont=NotoSansCJKsc-Regular,
  BoldItalicFont=NotoSansCJKsc-Bold
]{NotoSansCJKsc-Regular}
\setCJKfamilyfont{zhfs}[
  Path=fonts/,
  Extension=.otf,
  UprightFont=NotoSerifSC-Regular,
  BoldFont=NotoSerifSC-Bold,
  ItalicFont=NotoSerifSC-Regular,
  BoldItalicFont=NotoSerifSC-Bold
]{NotoSerifSC-Regular}
\setCJKfamilyfont{zhkai}[
  Path=fonts/,
  Extension=.otf,
  UprightFont=NotoSerifSC-Regular,
  BoldFont=NotoSerifSC-Bold,
  ItalicFont=NotoSerifSC-Regular,
  BoldItalicFont=NotoSerifSC-Bold
]{NotoSerifSC-Regular}
\setCJKfamilyfont{zhli}[
  Path=fonts/,
  Extension=.otf,
  UprightFont=NotoSerifSC-Regular,
  BoldFont=NotoSerifSC-Bold,
  ItalicFont=NotoSerifSC-Regular,
  BoldItalicFont=NotoSerifSC-Bold
]{NotoSerifSC-Regular}
\setCJKfamilyfont{zhyou}[
  Path=fonts/,
  Extension=.otf,
  UprightFont=NotoSansCJKsc-Regular,
  BoldFont=NotoSansCJKsc-Bold,
  ItalicFont=NotoSansCJKsc-Regular,
  BoldItalicFont=NotoSansCJKsc-Bold
]{NotoSansCJKsc-Regular}
\setCJKfamilyfont{zhmono}[
  Path=fonts/,
  Extension=.otf,
  UprightFont=NotoSansMonoCJKsc-Regular,
  BoldFont=NotoSansMonoCJKsc-Bold,
  ItalicFont=NotoSansMonoCJKsc-Regular,
  BoldItalicFont=NotoSansMonoCJKsc-Bold
]{NotoSansMonoCJKsc-Regular}

\renewcommand{\CJKrmdefault}{zhsong}
\renewcommand{\CJKsfdefault}{zhhei}
\renewcommand{\CJKttdefault}{zhmono}
\providecommand{\songti}{}
\renewcommand{\songti}{\CJKfamily{zhsong}}
\providecommand{\heiti}{}
\renewcommand{\heiti}{\CJKfamily{zhhei}}
\providecommand{\fangsong}{}
\renewcommand{\fangsong}{\CJKfamily{zhfs}}
\providecommand{\kaishu}{}
\renewcommand{\kaishu}{\CJKfamily{zhkai}}
\providecommand{\lishu}{}
\renewcommand{\lishu}{\CJKfamily{zhli}}
\providecommand{\youyuan}{}
\renewcommand{\youyuan}{\CJKfamily{zhyou}}

% Material Design Icons：图标字体作为项目资源随稿件分发，码点按字体清单登记。
% 图标宏使用固定码点和独立字号，避免把正文字体切换带进正文段落。
\newfontfamily\MaterialDesignIcons[
  Path=fonts/,
  Extension=.ttf,
  UprightFont=*,
  Renderer=OpenType
]{materialdesignicons-webfont}

% 颜色取自参考教材的蓝色栏目系统，辅色只承担资料分类。
\definecolor{bookblue}{HTML}{1F78B4}
\definecolor{bookcyan}{HTML}{54C7D2}
\definecolor{bookink}{HTML}{262626}
\definecolor{bookgray}{HTML}{68727A}
\definecolor{bookline}{HTML}{C9D3D9}
\definecolor{booklightblue}{HTML}{EAF5FB}
\definecolor{booklightgreen}{HTML}{E9F3E4}
\definecolor{booklightyellow}{HTML}{FFF4C9}
\definecolor{booklightviolet}{HTML}{EEE8F6}
\definecolor{bookorange}{HTML}{E38A24}
\definecolor{bookgreen}{HTML}{5D9C74}
\definecolor{bookviolet}{HTML}{8064A2}

% 图标采用零高度占位，字号和光学偏移可以按字面单独校准，不改变标题行高。
\newcommand{\mdiGlyphAt}[4]{%
  \raisebox{#4}[0pt][0pt]{%
    {\MaterialDesignIcons\fontsize{#1}{#1}\selectfont\color{#2}\char"#3}%
  }%
}
\newcommand{\mdiGlyph}[3][12]{\mdiGlyphAt{#1}{#2}{#3}{-0.14ex}}
\DeclareRobustCommand{\mdiBook}{\mdiGlyph[12]{bookblue}{F15D6}}
\DeclareRobustCommand{\mdiDocument}{\mdiGlyph[12]{bookgreen}{F09EE}}
% 问题图标默认使用头像轮廓；头像、聊天气泡和灯泡按字面分别校准。
\DeclareRobustCommand{\mdiQuestion}{\mdiGlyphAt{14.5}{bookorange}{F134B}{-0.42ex}}
\DeclareRobustCommand{\mdiHeadQuestion}{\mdiGlyphAt{14.5}{bookorange}{F134B}{-0.42ex}}
\DeclareRobustCommand{\mdiChatQuestion}{\mdiGlyphAt{14.5}{bookblue}{F1739}{-0.14ex}}
\DeclareRobustCommand{\mdiLightbulbQuestion}{\mdiGlyphAt{14.5}{bookorange}{F19E4}{-0.12ex}}
\DeclareRobustCommand{\mdiFileQuestion}{\mdiGlyphAt{13}{bookgreen}{F1036}{-0.10ex}}
\DeclareRobustCommand{\mdiTableQuestion}{\mdiGlyphAt{13}{bookgreen}{F1B21}{-0.08ex}}
\DeclareRobustCommand{\mdiBeakerQuestion}{\mdiGlyphAt{13}{bookviolet}{F1232}{-0.08ex}}
\DeclareRobustCommand{\mdiFolderQuestion}{\mdiGlyphAt{13}{bookblue}{F19CA}{-0.08ex}}
\DeclareRobustCommand{\mdiCommentQuestion}{\mdiGlyphAt{13}{bookblue}{F0186}{-0.06ex}}
\DeclareRobustCommand{\mdiMessageQuestion}{\mdiGlyphAt{13}{bookblue}{F173B}{-0.06ex}}
\DeclareRobustCommand{\mdiInfo}{\mdiGlyph[12]{bookviolet}{F02FD}}
\DeclareRobustCommand{\mdiQuote}{\mdiGlyph[12]{bookorange}{F1022}}
\DeclareRobustCommand{\mdiLink}{\mdiGlyph[11]{bookblue}{F0339}}
\DeclareRobustCommand{\mdiImage}{\mdiGlyph[11]{bookcyan}{F0976}}
\DeclareRobustCommand{\mdiTable}{\mdiGlyph[11]{bookgreen}{F04F1}}
\DeclareRobustCommand{\mdiKey}{\mdiGlyph[11]{bookorange}{F0336}}
\DeclareRobustCommand{\mdiSearch}{\mdiGlyph[11]{bookviolet}{F0349}}
\DeclareRobustCommand{\mdiCompass}{\mdiGlyph[11]{bookblue}{F018C}}
\DeclareRobustCommand{\mdiTimeline}{\mdiGlyph[11]{bookcyan}{F0BD2}}
\DeclareRobustCommand{\mdiBookmark}{\mdiGlyph[11]{bookorange}{F00C3}}
\DeclareRobustCommand{\mdiScale}{\mdiGlyph[11]{bookgreen}{F05D1}}
\DeclareRobustCommand{\mdiAtom}{\mdiGlyph[11]{bookviolet}{F0768}}
\DeclareRobustCommand{\mdiBrain}{\mdiGlyph[11]{bookblue}{F09D1}}
\DeclareRobustCommand{\mdiMap}{\mdiGlyph[11]{bookcyan}{F0982}}
\DeclareRobustCommand{\mdiTarget}{\mdiGlyph[11]{bookorange}{F04FE}}
\DeclareRobustCommand{\mdiClipboard}{\mdiGlyph[11]{bookgreen}{F0A38}}

\setstretch{1.20}
\setlength{\parindent}{2em}
\setlength{\parskip}{0pt}
\setlength{\emergencystretch}{2em}
% 数学行距与正文行距分开校准；行内公式不额外撑高中文行。
% equation/align 的编号按章递增，\eqref 可在正文中保留括号格式。
\mathsurround=0pt
\setlength{\abovedisplayskip}{8pt plus 2pt minus 2pt}
\setlength{\belowdisplayskip}{8pt plus 2pt minus 2pt}
\setlength{\abovedisplayshortskip}{4pt plus 1pt minus 1pt}
\setlength{\belowdisplayshortskip}{6pt plus 1pt minus 1pt}
\setlength{\jot}{4pt}
\numberwithin{equation}{chapter}
\renewcommand{\theequation}{\thechapter.\arabic{equation}}
\clubpenalty=9000
\widowpenalty=9000
\displaywidowpenalty=9000
\flushbottom
\setlist{nosep,leftmargin=2.2em}
\setlist[itemize]{label=\textcolor{bookblue}{\small\ensuremath{\diamond}}}

\hypersetup{
  unicode=true,
  colorlinks=true,
  linkcolor=bookblue,
  urlcolor=bookblue,
  citecolor=bookviolet,
  pdftitle={A4 中文书籍 LaTeX 模板},
  pdfauthor={}
}
\urlstyle{same}

% 正文出处的统一接口；默认放入脚注，也可以整体切换为文中括注。
% 可选参数用于在脚注中展开一个更完整的出处，正文仍保留短写。
\newif\ifbooksourceinline
\newcommand{\sourcecitesinline}{\booksourceinlinetrue}
\newcommand{\sourcecitesasfootnotes}{\booksourceinlinefalse}
\sourcecitesasfootnotes
\newcounter{booklastsourcefootnote}
\NewDocumentCommand{\sourcecite}{o m}{%
  \ifbooksourceinline
    （#2）%
  \else
    \IfValueTF{#1}{\footnote{#1}}{\footnote{#2}}%
    \setcounter{booklastsourcefootnote}{\value{footnote}}%
  \fi
}
\NewDocumentCommand{\sourceciteagain}{m}{%
  \ifbooksourceinline
    （#1）%
  \else
    \footnotemark[\value{booklastsourcefootnote}]%
  \fi
}

% 页眉只做导航，章首页清空页眉；页码放在外侧，适合装订后的双页阅读。
\pagestyle{fancy}
\fancyhf{}
\fancyhead[LE]{\small\color{bookgray}\leftmark}
\fancyhead[RO]{\small\color{bookgray}\rightmark}
\fancyfoot[LE]{\small\color{bookgray}\thepage}
\fancyfoot[RO]{\small\color{bookgray}\thepage}
\renewcommand{\headrulewidth}{0.35pt}
\renewcommand{\footrulewidth}{0pt}
\makeatletter
\renewcommand{\cleardoublepage}{\clearpage}
\renewcommand{\chaptermark}[1]{\markboth{第\thechapter 章\hspace{0.6em}#1}{}}
\renewcommand{\sectionmark}[1]{\markright{第\thesection 课\hspace{0.6em}#1}}
\makeatother
\fancypagestyle{chapterpage}{%
  \fancyhf{}
  \fancyfoot[LE]{\small\color{bookgray}\thepage}
  \fancyfoot[RO]{\small\color{bookgray}\thepage}
  \renewcommand{\headrulewidth}{0pt}
}
\fancypagestyle{plain}{%
  \fancyhf{}
  \fancyfoot[LE]{\small\color{bookgray}\thepage}
  \fancyfoot[RO]{\small\color{bookgray}\thepage}
  \renewcommand{\headrulewidth}{0pt}
}
\fancypagestyle{bookfrontpage}{%
  \fancyhf{}
  \fancyfoot[LE]{\small\color{bookgray}\thepage}
  \fancyfoot[RO]{\small\color{bookgray}\thepage}
  \renewcommand{\headrulewidth}{0pt}
}

% 样张用两级导航单位展示标题关系；正式书稿可以把这些接口映射到自己的目录层级。
\titleformat{\chapter}[block]
  {\normalfont\bfseries\color{bookink}\raggedright}
  {\raisebox{0.16\height}{\colorbox{bookblue}{\parbox[c][27mm][c]{35mm}{\centering\color{white}\fontsize{15}{19}\selectfont 第\thechapter 章}}}}
  {9mm}
  {\fontsize{34}{41}\selectfont}
\titlespacing*{\chapter}{0pt}{0pt}{16mm}

% 课是章节内部的第二级标题；课内标题使用 subsection。
\makeatletter
\@addtoreset{section}{chapter}
\makeatother

% 前言标题与正文从同一页开始；正文文件只需调用 \bookforeword。
\makeatletter
\newcommand{\bookforeword}{%
  \clearpage
  \phantomsection
  \addcontentsline{toc}{chapter}{前言}%
  \markboth{前言}{}%
  \thispagestyle{bookfrontpage}%
  \vspace*{16mm}%
  \noindent{\color{bookink}\fontsize{22}{28}\selectfont\bfseries 前言\par}%
  \vspace{6mm}%
}
\makeatother
\renewcommand{\thesection}{\arabic{section}}
\renewcommand{\thesubsection}{\arabic{section}.\arabic{subsection}}
\titleformat{\section}
  {\normalfont\mdseries\color{bookblue}\fontsize{20}{26}\selectfont\raggedright}
  {\raisebox{0.10\height}{\fcolorbox{bookblue}{booklightblue}{\parbox[c][12mm][c]{20mm}{\centering\color{bookblue}\fontsize{9.5}{12}\selectfont 第\thesection 课}}}}
  {6mm}
  {}
\titlespacing*{\section}{0pt}{8mm}{6mm}
\titleformat{\subsection}
  {\normalfont\bfseries\color{bookblue}\fontsize{14}{20}\selectfont}
  {\textcolor{bookblue}{\small\ensuremath{\diamond}}}
  {0.55em}
  {}
\titlespacing*{\subsection}{0pt}{5mm}{2mm}
\titleformat{\subsubsection}
  {\normalfont\bfseries\color{bookink}\fontsize{12.5}{18}\selectfont}
  {\thesubsection}
  {0.5em}
  {}
\titlespacing*{\subsubsection}{0pt}{4mm}{2mm}

% 目录层级：章名、课名和课内标题保持轻量秩序。
% 各 makebox 是编号栏，不是装饰空白；宽度按单/双位课号预留。
\setcounter{tocdepth}{2}
\titlecontents{chapter}[0.35em]
  {\addvspace{0.55em}\bfseries\color{bookink}\fontsize{15}{19}\selectfont}
  {\makebox[4.0em][l]{\thecontentslabel\hfill}}
  {}
  {\titlerule*[0.6pc]{.}\contentspage}
\titlecontents{section}[4.2em]
  {\color{bookblue}\fontsize{11.5}{16}\selectfont}
  {\makebox[4.8em][l]{第\thecontentslabel 课\hfill}}
  {}
  {\titlerule*[0.6pc]{.}\contentspage}
\titlecontents{subsection}[7.2em]
  {\color{bookgray}\fontsize{11}{15.5}\selectfont}
  {\makebox[3.2em][l]{\thecontentslabel\hfill}}
  {}
  {\titlerule*[0.6pc]{.}\contentspage}

% 目录使用独立页框，内容少时仍保持页面的整体比例。
\makeatletter
\newcommand{\bookcontents}{%
  \clearpage
  \thispagestyle{plain}%
  \begin{tikzpicture}[remember picture,overlay]
    \fill[booklightblue]
      ([xshift=10mm,yshift=-10mm]current page.north west)
      rectangle
      ([xshift=-10mm,yshift=10mm]current page.south east);
  \end{tikzpicture}%
  \vspace*{17mm}%
  \begin{center}
    {\color{bookink}\fontsize{25}{32}\selectfont\bfseries 目\hspace{0.7em}录}
  \end{center}
  \vspace{8mm}
  \@starttoc{toc}%
  \clearpage
}
\makeatother

% 栏目框：浅底色、细侧线、可见标题。正文不再被厚边框切成卡片。
% 四种颜色只标记栏目职责，标题和正文的字号仍由同一套比例控制。
\tcbset{
  bookbox/.style={
    enhanced,
    breakable,
    boxrule=0pt,
    frame hidden,
    sharp corners,
    left=5mm,
    right=4mm,
    top=3mm,
    bottom=3mm,
    before skip=5mm,
    after skip=6mm,
    fontupper=\normalsize,
    coltitle=bookink,
    fonttitle=\bfseries\normalsize,
    toptitle=1.6mm,
    bottomtitle=1.6mm,
    lefttitle=4mm,
    titlerule=0pt
  }
}
\newtcolorbox{readingbox}[1][阅读与思考]{
  bookbox,
  colback=booklightblue!62!white,
  borderline west={1.1pt}{0pt}{bookblue},
  colbacktitle=booklightblue,
  title={\mdiBook\hspace{0.8em}#1}
}
\newtcolorbox{sourcebox}[1][原典与材料]{
  bookbox,
  colback=booklightgreen!62!white,
  borderline west={1.1pt}{0pt}{bookgreen},
  colbacktitle=booklightgreen,
  title={\mdiDocument\hspace{0.8em}#1}
}
\newtcolorbox{questionbox}[1][问题]{
  bookbox,
  colback=booklightyellow!64!white,
  borderline west={1.1pt}{0pt}{bookorange},
  colbacktitle=booklightyellow,
  title={\mdiQuestion\hspace{0.8em}#1}
}
\newtcolorbox{noteBox}[1][提示]{
  bookbox,
  colback=booklightviolet!64!white,
  borderline west={1.1pt}{0pt}{bookviolet},
  colbacktitle=booklightviolet,
  title={\mdiInfo\hspace{0.8em}#1}
}
\makeatletter
\newcommand{\bookwrapclear}{%
  \par
  \ifnum\c@WF@wrappedlines>\z@
    \@tempcnta=\c@WF@wrappedlines
    \advance\@tempcnta by -1
    \vspace*{\dimexpr\@tempcnta\baselineskip\relax}%
  \fi
  \WFclear
}
\makeatother
% 引文折角固定为 5 mm；左下白色裁角和右上浅色折面共同构成纸页折角。
\newcommand{\quotefoldsize}{5mm}
\NewDocumentEnvironment{quotenote}{O{r} O{0.24\linewidth} m m}{%
  \begin{wrapfigure}{#1}{#2}
    \vspace{-2mm}
    \begin{tcolorbox}[
      enhanced,
      boxrule=0pt,
      frame hidden,
      sharp corners,
      width=\linewidth,
      colback=booklightyellow!82!white,
      left=3mm,
      right=3mm,
      top=2.5mm,
      bottom=6mm,
      before skip=0pt,
      after skip=0pt,
      overlay unbroken={%
        \fill[white] (frame.south west) -- ++(0,\quotefoldsize) -- ++(\quotefoldsize,-\quotefoldsize) -- cycle;
        \fill[bookorange!18!booklightyellow] ([yshift=\quotefoldsize]frame.south west) -- ++(\quotefoldsize,0) -- ++(0,-\quotefoldsize) -- cycle;
        \draw[bookorange!32!booklightyellow,line width=0.2pt] ([yshift=\quotefoldsize]frame.south west) -- ([xshift=\quotefoldsize]frame.south west);
      }
    ]
    \raggedright
    {\footnotesize\color{bookgray}\mdiQuote\hspace{0.45em}\textbf{#3}\par\scriptsize #4}\par\smallskip
    \fontsize{10.5}{15.5}\selectfont\itshape
}{%
  \end{tcolorbox}
  \end{wrapfigure}
}
% 小引文与侧边图像共用同一条 wrapfigure 路径，保证正文从盒子旁边开始排。
% sidequote 只能从段首调用；bookwrapclear 负责结束短段落留下的剩余环绕高度。
\newcommand{\bookquotecontent}[3]{%
  \begin{tcolorbox}[
    enhanced,
    boxrule=0pt,
    frame hidden,
    sharp corners,
    width=\linewidth,
    colback=booklightyellow!82!white,
    left=3mm,
    right=3mm,
    top=2.5mm,
    bottom=6mm,
    before skip=0pt,
    after skip=0pt,
    overlay unbroken={%
      \fill[white] (frame.south west) -- ++(0,\quotefoldsize) -- ++(\quotefoldsize,-\quotefoldsize) -- cycle;
      \fill[bookorange!18!booklightyellow] ([yshift=\quotefoldsize]frame.south west) -- ++(\quotefoldsize,0) -- ++(0,-\quotefoldsize) -- cycle;
      \draw[bookorange!32!booklightyellow,line width=0.2pt] ([yshift=\quotefoldsize]frame.south west) -- ([xshift=\quotefoldsize]frame.south west);
    }
  ]
  \raggedright
  {\footnotesize\color{bookgray}\mdiQuote\hspace{0.45em}\textbf{#1}\par\scriptsize #2}\par\smallskip
  \fontsize{10.5}{15.5}\selectfont\itshape #3
  \end{tcolorbox}%
}
\NewDocumentCommand{\sidequote}{O{r} O{0.24\linewidth} m m m}{%
  \ifvmode\Needspace{8\baselineskip}\fi
  \begin{wrapfigure}{#1}{#2}
    \vspace{-2mm}
    \bookquotecontent{#3}{#4}{#5}
  \end{wrapfigure}%
}

% fullquote 是自动引文接口：先按 40% 版心量高，短引文贴外侧，
% 放不下、会跨页或后文不足以环绕时自动改为 90% 版心的宽幅引文。
% 位置和环绕行数写入 aux，下一遍编译据此稳定排法；可选参数保留作兼容位。
\newcommand{\quotefoldoverlay}{%
  \fill[white] (frame.south west) -- ++(0,\quotefoldsize) -- ++(\quotefoldsize,-\quotefoldsize) -- cycle;
  \fill[bookorange!18!booklightyellow] ([yshift=\quotefoldsize]frame.south west) -- ++(\quotefoldsize,0) -- ++(0,-\quotefoldsize) -- cycle;
  \draw[bookorange!32!booklightyellow,line width=0.2pt] ([yshift=\quotefoldsize]frame.south west) -- ([xshift=\quotefoldsize]frame.south west);
}
\makeatletter
\def\fullquote@sidemode{side}
\def\gfullquote@mode{full}
\def\gfullquote@author{}
\def\gfullquote@source{}
\def\gfullquote@body{}
\newbox\fq@sbox
\newdimen\fq@w
\newdimen\fq@H
\newdimen\fq@Hgrid
\newdimen\fq@bottom
\newdimen\fq@asc
\newdimen\fq@dsc
\newcount\fq@lines
\newcount\fq@tmpcnt
\newcount\fq@est
\newcount\fq@pars
\newif\iffq@wrapping
\newif\iffq@counting
\newif\iffq@ok
\def\fq@key{}
\def\fq@prevkey{}
\def\fq@prevmode{}
\def\fq@place{o}
% 上一遍的引文记录：位置、页面、脚注区高度和是否锁定为 full。
\def\fq@rec#1#2#3{\expandafter\xdef\csname fq@#1@#2\endcsname{#3}}
\def\fq@foot#1#2{\expandafter\xdef\csname fq@foot@#1\endcsname{#2}}
\def\fq@val#1{\ifcsname fq@\fq@key @#1\endcsname\csname fq@\fq@key @#1\endcsname\else 0\fi}
\def\fq@textbottom{\dimexpr\paperheight-1in-\voffset-\topmargin-\headheight-\headsep-\textheight\relax}
\def\fq@foottop#1{\numexpr\fq@textbottom+\ifcsname fq@foot@#1\endcsname\csname fq@foot@#1\endcsname\else 0\fi\relax}
% 记录当前位置的纵坐标、绝对页号和页码；值在 shipout 时展开。
\def\fq@pos#1{%
  \pdfsavepos
  \edef\fq@tmp{\write\@auxout{%
    \string\fq@rec{\fq@key}{#1y}{\noexpand\the\pdflastypos}%
    \string\fq@rec{\fq@key}{#1p}{\noexpand\the\ReadonlyShipoutCounter}%
    \string\fq@rec{\fq@key}{#1g}{\noexpand\the\c@page}}}%
  \fq@tmp
  \ifvmode\if@nobreak\nobreak\fi\fi}
\pretocmd{\@makecol}{%
  \ifvoid\footins\else
    \immediate\write\@auxout{\string\fq@foot{\the\numexpr\ReadonlyShipoutCounter+1\relax}{\the\numexpr\dimexpr\ht\footins+\dp\footins+\skip\footins\relax\relax}}%
  \fi}{}{}
\AddToHook{para/begin}{\iffq@counting\global\advance\fq@pars\@ne\fi}
% wrapfig 的默认余量检查会把已经量好的 side 盒子下沉；自动引文使用已记录的位置直接放置。
\let\fq@orig@putfigmaybe\WF@putfigmaybe
\def\WF@putfigmaybe{%
  \iffq@wrapping
    \global\fq@wrappingfalse
    \expandafter\fq@putfig
  \else
    \expandafter\fq@orig@putfigmaybe
  \fi}
\def\fq@putfig{%
  \vskip-\parskip
  \noindent
  \global\WF@floatfalse
  {\ifodd\if@twoside\c@page\else\@ne\fi
    \lccode`i`l\lccode`o`r\else \lccode`i`r\lccode`o`l\fi
    \xdef\WF@place{\the\lccode\lccode\WF@place}}%
  \hbox to\z@{%
   \@tempdima\wd\WF@box \@tempdimb\WF@ovh
   \advance\@tempdima-\@tempdimb \advance\@tempdima\columnsep
   \advance\@tempdimb\hsize \advance\@tempdimb-\@tempdima
   \xdef\WF@adjlw{\the\@tempdima}%
   \ifnum `l=\WF@place
    \hss
    \def\@tempa{\kern\columnsep}%
   \else
    \@tempdima\z@
    \kern\@tempdimb \kern\columnsep
    \def\@tempa{\hss}%
   \fi
   \ifdim\@tempdimb<\hsize
    \xdef\WF@wrapil{\the\@tempdima \the\@tempdimb}%
    \xdef\WF@adjtlm{\the\@tempdima}%
   \else
    \xdef\WF@wrapil{\z@ \the\hsize}%
    \xdef\WF@adjlw{\z@}\xdef\WF@adjtlm{\z@}%
   \fi
   \dp\WF@box\z@ \box\WF@box \@tempa
  }%
  \aftergroup\WF@startwrapping}
\def\fq@setplace#1{\ifodd#1\relax\gdef\fq@place{r}\else\gdef\fq@place{l}\fi}
% 上一遍排成 side：盒底不压脚注、上下在同一页，且后文把盒子环绕完。
\def\fq@checkside{%
  \fq@oktrue
  \ifnum\fq@val{short}>\z@ \fq@okfalse\fi
  \ifnum\fq@val{tp}=\fq@val{bp}\relax\else\fq@okfalse\fi
  \ifnum\numexpr\fq@val{by}-\dimexpr2pt\relax\relax<\fq@foottop{\fq@val{bp}}\relax\fq@okfalse\fi
  \iffq@ok
    \fq@setplace{\fq@val{tg}}%
  \else
    \gdef\gfullquote@mode{full}%
    \immediate\write\@auxout{\string\fq@rec{\fq@key}{lock}{\fq@val{tg}}}%
  \fi}
% 上一遍排成 full：按 full 起点估算 side 盒的位置，并按后文行数估算能否环绕完。
\def\fq@checkfull{%
  \fq@okfalse
  \ifnum\numexpr\fq@val{rg}-\fq@val{lock}\relax<2 \ifnum\fq@val{rg}<\fq@val{lock}\relax\fq@okfalse\else\fq@oktrue\fi\else\fq@okfalse\fi
  \iffq@ok
    \fq@okfalse
    \immediate\write\@auxout{\string\fq@rec{\fq@key}{lock}{\fq@val{lock}}}%
  \else\ifcsname fq@\fq@key @ry\endcsname
    \fq@oktrue
    \ifnum\numexpr\fq@val{ry}-\dimexpr\baselineskip-\fq@asc-\fq@dsc+\fq@Hgrid+2pt\relax\relax<\fq@foottop{\fq@val{rp}}\relax
      \fq@okfalse
    \fi
    \ifcsname fq@\fq@key @cy\endcsname
      \ifnum\fq@val{cp}=\fq@val{ap}\relax
        \fq@tmpcnt=\numexpr(\fq@val{ay}-\fq@val{cy})/\dimexpr\baselineskip\relax\relax
        \advance\fq@tmpcnt-\fq@val{np}\relax
        \fq@est=\numexpr\fq@tmpcnt*\dimexpr\linewidth\relax/\dimexpr\linewidth-\fq@w-\columnsep\relax+\fq@val{np}-1\relax
        \ifnum\fq@est<\fq@lines\fq@okfalse\fi
      \fi
    \fi
  \fi\fi
  \iffq@ok\fq@setplace{\fq@val{rg}}\else\gdef\gfullquote@mode{full}\fi}
\def\fq@decide{%
  \gdef\gfullquote@mode{side}\gdef\fq@place{o}%
  \ifcsname fq@\fq@key @mode\endcsname
    \if s\csname fq@\fq@key @mode\endcsname
      \fq@checkside
    \else
      \fq@checkfull
    \fi
  \fi}
\newcommand{\fq@sidetcb}[1]{%
  \begin{tcolorbox}[
    enhanced,
    boxrule=0pt,
    frame hidden,
    sharp corners,
    width=\fq@w,
    colback=booklightyellow!82!white,
    left=3mm,right=3mm,top=2.5mm,bottom=6mm,
    before skip=0pt,after skip=0pt,
    overlay unbroken={\quotefoldoverlay},
    #1
  ]
  \raggedright
  {\interlinepenalty=10000\footnotesize\color{bookgray}\mdiQuote\hspace{0.45em}\textbf{\gfullquote@author}\par\scriptsize \gfullquote@source\par}\nobreak\smallskip\nobreak
  \normalsize\itshape \gfullquote@body
  \end{tcolorbox}}
\newcommand{\fullquotewrapclear}{%
  \par
  \ifx\fq@prevkey\@empty\else
    \begingroup
    \let\fq@key\fq@prevkey
    \ifx\fq@prevmode\fullquote@sidemode
      \ifvoid\WF@box
        \ifnum\c@WF@wrappedlines>\@ne
          \immediate\write\@auxout{\string\fq@rec{\fq@key}{short}{\the\numexpr\c@WF@wrappedlines-1\relax}}%
        \fi
      \else
        \immediate\write\@auxout{\string\fq@rec{\fq@key}{short}{99}}%
      \fi
    \else
      \fq@pos{c}%
      \immediate\write\@auxout{\string\fq@rec{\fq@key}{np}{\the\fq@pars}}%
    \fi
    \endgroup
    \global\let\fq@prevkey\@empty
    \global\fq@countingfalse
  \fi
  \global\fq@wrappingfalse
  \ifnum\c@WF@wrappedlines>\@ne
    \@tempcnta=\c@WF@wrappedlines
    \advance\@tempcnta by -1
    \vspace*{\dimexpr\@tempcnta\baselineskip\relax}%
  \fi
  \WFclear
}
\newcommand{\fullquote@after}{%
  \par
  \global\let\fq@prevkey\fq@key
  \global\let\fq@prevmode\gfullquote@mode
  \ifx\gfullquote@mode\fullquote@sidemode
    \immediate\write\@auxout{\string\fq@rec{\fq@key}{mode}{s}}%
    \begin{wrapfigure}[\the\fq@lines]{\fq@place}{\fq@w}
      \hrule\@height\fq@asc\@depth\z@\@width\z@
      \vskip-\fq@asc
      \fq@pos{t}%
      \setbox\fq@sbox\vbox{\hsize\fq@w\fq@sidetcb{height=\fq@Hgrid,bottom=\fq@bottom}}%
      \box\fq@sbox
      \fq@pos{b}%
    \end{wrapfigure}%
    \global\fq@wrappingtrue
  \else
    \immediate\write\@auxout{\string\fq@rec{\fq@key}{mode}{f}}%
    \fq@pos{r}%
    \begin{tcolorbox}[
      enhanced,
      breakable,
      boxrule=0pt,
      frame hidden,
      sharp corners,
      width=0.90\linewidth,
      center,
      lines before break=2,
      colback=booklightyellow!82!white,
      left=5mm,right=5mm,top=4mm,bottom=6mm,
      before skip=5mm,after skip=6mm,
      overlay unbroken={\quotefoldoverlay},
      overlay last={\quotefoldoverlay}
    ]
    \raggedright
    {\interlinepenalty=10000\footnotesize\color{bookgray}\mdiQuote\hspace{0.55em}\textbf{\gfullquote@author}\enspace \gfullquote@source\par}\nobreak\smallskip\nobreak
    \normalsize\itshape \gfullquote@body
    \end{tcolorbox}
    \par
    \fq@pos{a}%
    \global\fq@pars\z@
    \global\fq@countingtrue
  \fi
}
\NewDocumentEnvironment{fullquote}{O{} m m +b}{%
  \fullquotewrapclear
  \gdef\gfullquote@author{#2}%
  \gdef\gfullquote@source{#3}%
  \gdef\gfullquote@body{#4}%
  \xdef\fq@key{\mdfivesum{\detokenize{#2}\detokenize{#3}\detokenize{#4}}}%
  \gdef\gfullquote@mode{full}%
  \global\fq@w=0.40\linewidth
  \global\let\fq@savednobreak\if@nobreak
  \setbox\fq@sbox\vbox{\hsize\fq@w\fq@sidetcb{}}%
  \global\let\if@nobreak\fq@savednobreak
  \global\fq@H=\dimexpr\ht\fq@sbox+\dp\fq@sbox\relax
  \ifdim\fq@H<1.6\fq@w
    \setbox\fq@sbox\hbox{\normalfont\normalsize 国}%
    \global\fq@asc=\ht\fq@sbox
    \global\fq@dsc=\dp\fq@sbox
    \fq@tmpcnt=\numexpr\dimexpr\fq@H-\fq@asc-\fq@dsc-0.3\baselineskip+\baselineskip\relax-1\relax
    \divide\fq@tmpcnt by \numexpr\dimexpr\baselineskip\relax\relax
    \advance\fq@tmpcnt\@ne
    \ifnum\fq@tmpcnt<\@ne\fq@tmpcnt\@ne\fi
    \global\fq@lines=\fq@tmpcnt
    \global\fq@Hgrid=\dimexpr\fq@asc+\baselineskip*(\fq@lines-1)+\fq@dsc\relax
    \global\fq@bottom=6mm
    \ifdim\fq@H>\fq@Hgrid
      \global\advance\fq@bottom by -\dimexpr\fq@H-\fq@Hgrid\relax
    \fi
    \fq@decide
  \fi
}{ }
\AfterEndEnvironment{fullquote}{\fullquote@after}
\let\bookwrapclear\fullquotewrapclear
\makeatother

% 表格统一处理前后呼吸，避免表尾直接撞上下一段正文。
% 采用 minipage 而非 center，消除 center 环境自带的额外 topsep。
\NewDocumentEnvironment{booktable}{O{0.92\linewidth} m +b}{%
  \par\addvspace{1mm}%
  \noindent
  \begin{minipage}{\linewidth}
    \centering
    \small
    \renewcommand{\arraystretch}{1.18}%
    \begin{tabularx}{#1}{#2}
      #3
    \end{tabularx}
  \end{minipage}
  \par\addvspace{6mm}%
}{}
\newcommand{\tablecaption}[1]{%
  \par\addvspace{2mm}\noindent
  {\footnotesize\color{bookgray}\mdiTable\hspace{0.55em}#1}\par\vspace{1mm}
}

% 参考书中的“相关链接”是页面节奏，而不是一只大盒子。
\newcommand{\relatedlink}[1]{%
  \par\Needspace{4\baselineskip}%
  \vspace{2mm}\noindent
  {\color{bookblue}\mdiLink\hspace{0.65em}\bfseries 相关链接}\hspace{0.6em}%
  {\color{bookline}\leaders\hbox{\rule[0.45ex]{1.1mm}{0.35pt}}\hfill\kern0pt}\par
  {\small\color{bookgray}#1\par}\vspace{2mm}
}

% 章导言和图像位置的统一接口。样张中用抽象占位块，正式稿替换为真实图像。
% 占位块只表示网格尺寸；正式图像由 fullimage/sideimage 按原比例载入。
\newcommand{\chapterlead}[1]{%
  \begin{tcolorbox}[
    enhanced,
    frame hidden,
    boxrule=0pt,
    colback=white,
    borderline west={1.2pt}{0pt}{bookcyan},
    left=5mm,
    right=0pt,
    top=0pt,
    bottom=0pt,
    before skip=4mm,
    after skip=5mm
  ]
  \normalsize\color{bookblue}#1
  \end{tcolorbox}
}
\newcommand{\lessonlead}[1]{%
  \begin{tcolorbox}[
    enhanced,
    frame hidden,
    boxrule=0pt,
    colback=white,
    borderline west={1.2pt}{0pt}{bookcyan},
    left=5mm,
    right=0pt,
    top=0pt,
    bottom=0pt,
    before skip=3mm,
    after skip=5mm
  ]
  \normalsize\color{bookblue}#1
  \end{tcolorbox}
}
\newcommand{\chaptervisual}[1]{%
  \begin{center}
    \begin{tikzpicture}
      \node[draw=bookline,fill=booklightblue,minimum width=0.82\linewidth,minimum height=120mm,align=center,text=bookblue,font=\small,rounded corners=2mm] (cv) {#1};
      \draw[bookcyan,line width=1.1pt] (cv.north west) -- (cv.south east);
      \draw[bookcyan,line width=1.1pt] (cv.south west) -- (cv.north east);
    \end{tikzpicture}
  \end{center}
}
\newcommand{\chapterabstract}[1]{%
  \begin{tcolorbox}[
    enhanced,
    frame hidden,
    boxrule=0pt,
    colback=white,
    borderline west={1.2pt}{0pt}{bookgray},
    left=5mm,
    right=0pt,
    top=0pt,
    bottom=0pt,
    before skip=4mm,
    after skip=5mm
  ]
  \normalsize\color{bookgray}#1
  \end{tcolorbox}
}
\newcommand{\fullplaceholder}[2][0.48\textwidth]{%
  \begin{center}
    \begin{tikzpicture}
      \node[draw=bookline,fill=booklightblue,minimum width=#1,minimum height=42mm,align=center,text=bookblue,font=\small] (img) {#2};
      \draw[bookcyan,line width=1.2pt] (img.north west) -- (img.south east);
      \draw[bookcyan,line width=1.2pt] (img.south west) -- (img.north east);
    \end{tikzpicture}
  \end{center}
}
\newcommand{\imageplaceholder}{\fullplaceholder}
\newcommand{\fullimage}[2][\linewidth]{%
  \begin{center}
    \includegraphics[width=#1]{#2}
  \end{center}
}
\newcommand{\sideplaceholder}[4][r]{%
  \Needspace{8\baselineskip}%
  \begin{wrapfigure}{#1}{#2}
    \vspace{-2mm}
    \centering
    {\setlength{\fboxsep}{0pt}%
      \fcolorbox{bookline}{booklightblue}{%
        \parbox[c][#3][c]{\dimexpr\linewidth-2\fboxrule\relax}{%
          \centering\small\color{bookblue}图像位置
        }%
      }%
    }%
    \captionline{#4}
  \end{wrapfigure}
}
\newcommand{\sideimage}[4][r]{%
  \Needspace{8\baselineskip}%
  \begin{wrapfigure}{#1}{#2}
    \vspace{-2mm}
    \centering
    \includegraphics[width=\linewidth]{#3}
    \captionline{#4}
  \end{wrapfigure}
}
\newcommand{\captionline}[1]{\begin{center}\footnotesize\color{bookgray}\mdiImage\hspace{0.45em}#1\end{center}}
\newcommand{\term}[1]{\textbf{#1}}
\newcommand{\foreignterm}[1]{\textit{#1}}

% 让章节首面使用专门的页眉状态。
\makeatletter
\pretocmd{\chapter}{\fullquotewrapclear\thispagestyle{chapterpage}}{}{}
\pretocmd{\section}{\fullquotewrapclear\Needspace{5\baselineskip}}{}{}
\pretocmd{\subsection}{\fullquotewrapclear\Needspace{4\baselineskip}}{}{}
\makeatother
